% Terjemahan Bahasa Indonesia dari prelude.tex baris 9-29.
% Sumber: Methods of Algebra, Volume 2, commit
% 9a5803ff77dd3257484cb177f851a73770a59dd3 (CC BY 4.0).
% stable-unit-id: o014.aljabr2.prelude.linear-algebra

\chapter*{Pendahuluan}
\hypertarget{o014.aljabr2.prelude}{ }

% segment-id: o014.li2.prelude.p001
Karya yang disajikan kepada pembaca ini merupakan kelanjutan dari
\cite{Li1} (selanjutnya disebut Jilid I). Tujuannya adalah memperkenalkan
serangkaian metode, gagasan, dan teknik yang secara wajar dapat dihimpun
di bawah nama aljabar linear, dengan mengandaikan bahwa pembaca telah
mengenal struktur-struktur dasar dalam aljabar. Metode-metode ini dipakai
di seluruh matematika kontemporer. Agar dapat menjawab kebutuhan praktis
sekomprehensif mungkin, teknik-teknik terkait perlu ditempa menjadi bentuk yang
lebih murni dan lebih ringkas. Kategori dan funktor merupakan bahasa yang
tak tergantikan untuk tujuan tersebut.

% segment-id: o014.li2.prelude.p002
Buku ini terbagi menjadi tiga bagian: bagian dalam, bagian luar,
dan lampiran. Sasaran utamanya ialah mahasiswa sarjana tingkat lanjut,
mahasiswa pascasarjana, peneliti, pengajar, dan pembelajar mandiri yang
menekuni atau berminat pada pokok-pokok ini. Pengetahuan prasyaratnya
mencakup grup, gelanggang, modul, lapangan, dan struktur aljabar lain, serta
teori kategori. Jilid I mencakup seluruh landasan tersebut, tetapi buku
teks bermutu lainnya juga memadai, meskipun mungkin terdapat sedikit
perbedaan notasi.
Motivasi penulisan telah dijelaskan dalam pendahuluan Jilid I; tujuannya
tetap sama dan tidak perlu diulang di sini.

\section*{Apakah Aljabar Linear Itu?}
\addcontentsline{toc}{section}{Apakah Aljabar Linear Itu?}
\hypertarget{o014.aljabr2.prelude.linear-algebra}{ }
\index{aljabar linear (linear algebra)}

% segment-id: o014.li2.prelude.linear-algebra.p001
Dalam pengertian yang lazim di kalangan ahli aljabar, kata ``linear''
secara luas merujuk pada sifat-sifat yang berhubungan dengan struktur
yang memiliki penjumlahan beserta operasi inversnya, yaitu pengurangan.
Penjumlahan dan pengurangan selanjutnya memberikan operasi perkalian
dengan bilangan bulat; generalisasi langsungnya adalah perkalian skalar
oleh unsur suatu gelanggang $R$. Contoh prototipikalnya ialah modul kiri
dan modul kanan atas $R$, sedangkan modul-$\Z$ tidak lain adalah grup
abelian. Untuk suatu modul-$R$ (mulai sekarang modul kiri, kecuali dinyatakan
lain), kita dapat mempelajari submodul dan modul kuosien (modul hasil bagi),
serta kernel
$\Ker(f)$, kokernel $\Coker(f)$, dan citra $\Image(f)$ dari suatu
homomorfisme modul $f$. Semuanya merupakan objek yang sudah dikenal dari
aljabar dasar; setidaknya, kasus ketika $R$ adalah lapangan, yakni kasus
ruang vektor, telah dikenal luas.

% segment-id: o014.li2.prelude.linear-algebra.p002
Sejak abad ke-20, praktik matematika secara bertahap memperlihatkan bahwa
memperluas kajian modul-$R$ dan homomorfismenya ke kompleks modul bukan
hanya bermanfaat dan memudahkan, melainkan sering kali diperlukan.
Kompleks semacam itu adalah deretan homomorfisme modul yang memenuhi
$d^{n+1}d^n=0$:
\[
  \cdots \longrightarrow X^{n-1} \xrightarrow{d^{n-1}} X^n
  \xrightarrow{d^n} X^{n+1} \xrightarrow{d^{n+1}} \cdots .
\]
Syarat pada homomorfisme tersebut kadang ditulis $d^2=0$, sedangkan
seluruh data kompleks sering disingkat menjadi $(X^n,d^n)_n$, $(X,d)$,
atau $X$. Karena digunakan indeks atas yang meningkat, kompleks ini juga
disebut kompleks kokrantai (\emph{cochain complex}). Jika digunakan indeks
bawah yang menurun,
$\cdots\to X_n\xrightarrow{d_n}X_{n-1}\to\cdots$, objek matematis yang
diperoleh disebut kompleks rantai; perbedaannya dengan kompleks di atas
hanya bersifat formal.

% segment-id: o014.li2.prelude.linear-algebra.p003
Suatu modul-$R$ tunggal $M$ dapat dipandang sebagai kasus khusus kompleks,
dengan $X^0=M$ dan semua suku lainnya nol. Untuk suatu kompleks $X$,
karena $d^nd^{n-1}=0$, kohomologi ke-$n$ dapat didefinisikan sebagai
modul kuosien
\[
  \Hm^n(X):=\Ker\left(d^n\right)\big/\Image\left(d^{n-1}\right).
\]

% segment-id: o014.li2.prelude.linear-algebra.p004
Kompleks yang memenuhi $\Hm^n(X)=0$ untuk setiap $n$ disebut barisan
eksak atau kompleks asiklik. Untuk kompleks rantai, secara bersesuaian
kita memiliki homologi
$\Hm_n(X):=\Ker(d_n)/\Image(d_{n+1})$. Kohomologi suatu kompleks (atau
homologi suatu kompleks rantai) sering menyimpan informasi penting
tentang objek matematis yang dikaji.

\enlargethispage{3\baselineskip}
% segment-id: o014.li2.prelude.linear-algebra.p005
Untuk modul atau struktur matematis lain yang memiliki operasi linear,
misalnya objek dalam kategori abelian yang akan diperkenalkan kemudian,
kajian kompleks beserta kohomologinya merupakan isi klasik bidang yang
disebut aljabar homologis. Bidang ini merupakan subhimpunan sejati dari
``aljabar linear'' dalam arti hakikinya dan sekaligus merupakan inti buku
ini. Riwayat singkatnya akan dibahas berikutnya.
